\documentclass[11pt,a4paper]{article}
\pagestyle{empty}
\usepackage[spanish]{babel}
\usepackage[utf8]{inputenc}
\usepackage{ragged2e}%justifica títulos
\usepackage{pdfpages}
\usepackage{xparse}
\DeclareGraphicsExtensions{.pdf,.png,.jpg}
\usepackage[rmargin=2.54cm,lmargin=2.54cm,top=2.54cm,bottom=2.54cm]{geometry}
\usepackage{amsmath}
\usepackage{amsfonts}
\usepackage{amssymb}
\usepackage{stackengine}
\usepackage{graphics}
\usepackage{graphicx}
\usepackage{tikz}
\usepackage{cancel}
\usepackage{color}
\usepackage{fancyhdr}
\pagestyle{fancy}
\usepackage{hyperref}
\usepackage{apacite}
\bibliographystyle{apacite}
\usepackage{pdfpages}
%%%%%%%%%%

\rhead{Jeff Preston Bezos}
\lhead{Amazon.com Inc. }
\fancyfoot[LE,RO]{\thepage}
\fancyfoot[CE,CO]{Comité Empresarial}
\bibliographystyle{apacite}

\title{Folder}
\date{12 de Abril de 2021}

\begin{document} %Aquí empieza

%\maketitle\thispagestyle{fancy}

\includepdf[pages=1]{portada}

\tableofcontents
\vspace{20cm}

\section{Información empresa}
\subsection{Historia general e indicadores}

\noindent  Creé Amazon en 1994 en Washington, Estados Unidos. Dado el enorme desarrollo que estaba teniedo internet, decidí emprender una \textit{startup }de venta de libros online con dinero que me otorgaron mis padres. Gracias a un trabajo arduo de nuestro equipo, crecimos exponencialmente, aumentando la diversidad de productos que los usuarios podían comprar y vender en nuestra plataforma. En 1995 tuvimos nuestra primera ronda de inversión, donde levantamos capital por más de 8 millones de dólares por parte del fondo de \textit{venture capital }Kleiner Perkins Caufield \& Byers. En 1997 decidimos convertirnos en una empresa pública. Hicimos nuestro IPO en el \textit{Nasdaq} cotizando a tan solo 18\$ la acción. Hoy, 24 años después, cada acción de nuestra compañía vale más de 3000\$. Este increíble éxito en bolsa hizo que en 2018 la capitalización bursátil de nuestra compañía sea de más de un billón de dólares. \cite{bbc}

\subsection{Finanzas}

\noindent Como empresa pública tenemos la responsabilidad de mostrar nuestros resultados financieros con transparencia a nuestros inversionistas. 

\vspace{0.5cm}

\includegraphics[scale=0.7]{balance.jpeg}


\noindent Los datos mostrados se obtuvieron de la bolsa de valores del \textit{Nasdaq}\footnote{Para más detalles consultar https://www.nasdaq.com/market-activity/stocks/amzn/financials}. Muestran nuestro flujo de caja en estos últimos 4 años. Como las cifras muestran, no pagamos dividendos hacia nuestros accionistas. Consideramos que los ingresos de Amazon deben destinarse de vuelta en la empresa, en los salarios de nuestros trabajadores y en I+D.

\subsection{Cultura}

\noindent Amazon tiene una cultura empresarial que buscan explotar todas las capacidades de nuestros trabajadores. Diseñamos 14 principios de liderazgo de los cuales no sentimos muy orgullosos \footnote{Para más detalles consultar \url{https://aws.amazon.com/es/careers/culture/}}:
\begin{itemize}
\item Fijación por el cliente
\item Asumir la responsabilidad
\item Inventar y simplificar
\item Tener razón muy seguido
\item Aprender y sentir curiosidad
\item Contratar y desarrollar a los mejores
\item Insistir en los estándares más altos
\item Pensar en grande
\item Tender a la acción
\item Ser austero
\item Ganarse la confianza
\item Ir más allá
\item Tener agallas, expresar desacuerdo y comprometerse
\item Lograr resultados
\end{itemize}

\noindent Además de esto, nuestros ejecutivos se caracterizan por su honestidad y su compromiso con nuestros clientes.

\subsection{Organismos a los que pertenecemos }

\noindent Somos parte de diversas comunidades empresariales que buscan fomentar del desarrollo de  tecnología y ciencia. Algunas de las más destacadas son:
\begin{itemize}
\item Alliance for Open Media
\item Internet Association
\item World Wide Web Consortium
\item  Asociación de Tarjetas SD
\item Zigbee Alliance
\end{itemize}

\subsection{Aliados}

\noindent Como varios empresarios han expuesto, en los negocios no existen amigos. No contamos con verdaderas alianzas con otras multinacionales. Eso sí, tenemos la fortuna de recibir el apoyo de diversos gobiernos alrededor del mundo donde operamos. Su colaboración logística (especialmente en temas de servidores de internet) nos ha ayudado a llegar a este punto. Nos sentimos especialmente agradecidos con la colaboración de EEUU, país donde hemos crecido. Si bien es cierto que hemos tenido algunas discrepancia con el congreso \cite{sundar} estamos dispuestos a ayudar en lo que sea necesario a proteger los datos de nuestros consumidores.

\subsection{Competidores}

\noindent Amazon nunca ha tenido intenciones de desmeritar o sabotear la labor de nuestros competidores. Sin embargo, hay que reconocer que buscamos sobresalir en mercados contra multinacionales como Alibaba en China y Mercadolibre en América Latina. Constantemente estamos mejorando la calidad de nuestro servicio para superarlos.

\section{Tema 1. Limitaciones actuales de las multinacionales durante su integración en Latinoamérica }

\subsection{Función del comité respecto al tema}

\noindent El comité empresarial, dado que integra tanto países como multinacionales, busca llegar a acuerdos entre empresas y naciones para acordar beneficios comunes. En este caso, se habla de minimizar los efectos negativos de las multinacionales en países latinoamericanos. Amazon buscará expandir su influencia en la región sin transgredir los derechos de los trabajadores. 

\subsection{Causas}

\noindent Las multinacionales con pocos escrúpulos, diferentes a Amazon, buscan evadir regulaciones que afectan el crecimiento de su empresa. Estas pueden ser no declarar impuestos a través de paraísos fiscales, hacer \textit{lobby} para que los políticos remuevan regulaciones ambientales o derechos laborales; o simplemente sobornar jueces que investigan sus crímenes.

La verdadera causa de esta situación se da por la gran corrupción de los políticos lationamericanos. Los países de la región suelen estar en el ranking de los más corruptos del planeta \cite{cor}. Existen diversas razones que conducen a este entorno de deshonestidad, pero las 2 mayores son la fragilidad de la instituciones encargadas de investigar a los políticos, y la gran cantidad de recursos de los partidos políticos tradicionales. ¿De dónde proviene el dinero de las campañas de estos partidos? De donaciones principalmente. Donaciones de empresas que buscan que sus candidatos lleguen al poder y así defiendan sus intereses.

\subsection{Consecuencias}

\noindent Los gobiernos lationoamericanos han perdido su legitimidad ante tantos escándalos y el derroche de recursos causado por la corrupción. Por esta razón, la ciudadanía suele tener una mala percepción de las multinacionales, ignorando el impacto positivo en la economía que generan empresas como Amazon. 
\subsection{Involucrados en el conflicto}

\subsubsection{Instituciones}

\noindent Existen varias multinacionales investigadas por estos abusos que están dentro de la comisión. Aquí una breve descripción de algunas:

\begin{itemize}
\item \textbf{HSBC:} Se declaró culpable al banco HSBC por lavar dinero del cartel de Sinaloa \cite{hsbc}.
\item \textbf{NTR Metals:} NTR Metals compraba oro de minerías ilegales —usualmente provenientes de Perú— para lavar dinero de narcotraficantes \cite{ntr}. 
\item \textbf{JP Morgan:} Este banco ha sido acusado en varias ocasiones de evasión de impuestos cuando hace operaciones en Latinoamérica, pero su papel más polémico fue su participación en el desarrollo de la crisis financiera de 2008; vendiendo CDO's\footnote{Collateralized Debt Obligation} con préstamos en \textit{default}.
\item \textbf{Odebrecht:} Esta famosa empresa sobornaba a políticos en diversas naciones de Latinoamérica para que les otorgaran licencias de construcción, mostrando una competencia desleal frente a otras firmas de construcción \cite{odb}.
\item \textbf{Formosa}: Esta empresa contaminó secretamente fuentes hídricas en Point Comfort, Texas y en otros países asiáticos \cite{for}.
\end{itemize}

\noindent Pasando al tema ambiental del segundo tema; absolutamente todas las empresa del comité que generan productos físicos —a diferencia de empresas de servicios como los bancos— contribuyen al proceso de extracción de recursos indirecta o directamente. Esto no implica que hayan infringido normas ambientales de los países en que operan, simplemente están acostumbrados al ritmo consumista del mercado. Por esto Amazon quiere hacer hincapié en que los problemas ambientales que estamos sufriendo no es cumpla directa de las empresas, sino de los consumidores. Ellos son quienes demandan productos sin pensar en el impacto que estos tendrán.

\subsubsection{Países}

\noindent Los países involucrados cuentan con empresas estatales que han generado efectos dañinos en el ambiente. EEUU y China son las naciones que más cantidad de $CO_{2}$ producen en el planeta \cite{co}. Por otra parte, Rusia tiene el monopolio de recursos naturales donde implementa estructuras como oleoductos de gas natural que han contaminado el entorno europeo \cite{gas}.

\subsection{Políticas de la delegación frente al tema}

\subsubsection{Antecedentes}

\noindent Amazon ha iniciado operaciones comerciales en grandes naciones lationoamericanas como México Brasil y Colombia; y se siente orgulloso de no verse involucrado en ningún escándalo legal. Reconocemos haber aportado dinero a campañas políticas en EEUU, aclarando que esto no hace parte de ningún crimen, ya que el \textit{lobby} es legal en este país. Sus aportes no representan ningún tipo de conducta inmoral, buscamos hacer alianzas con el gobierno de EEUU para iniciativas como la gestión de entregas de vacunas y desarrollo de inteligencia artificial \cite{lo}.

\subsubsection{Posibles soluciones}

\noindent Las propuestas de Amazon buscaran que las políticas que implementarán los estados fortalezcan su institucionalidad y evite acuerdos ilegales entre políticos y empresarios. Eso sí, estas políticas no deberían oponerse al libre mercado ni a la globalización. Denunciaremos con todo coraje a las delegaciones del comité que estén utilizando su influencia para intereses particulares. 

Para promover iniciativas que fomenten la transparencia  podemos utilizar alianzas con organismos multilaterales como la ONU o el Banco de Desarrollo de América Latina (CAF).

\subsection{Preguntas orientadoras}

\noindent La mayoría de preguntas orientadoras se pueden responder diciendo que tenemos los requisitos legales para operar en los países latinoamericanos que hemos mencionado. Hemos iniciado nuestra expansión en este continente para derrocar a nuestro competidor argentino Mercadolibre. 

Si bien hemos fomentado el empleo en nuestros almacenes de distribución, también hemos abierto puestos de trabajos para programadores y profesionales en atención al cliente. Consideramos que latinoamerica tiene un gran potencial de desarrollo tecnológico. 


\section{Tema 2. Influencia del consumismo dentro del uso de los recursos planetarios}

\subsection{Causas}

\noindent El sistema económico en el que vivimos ha fomentado el consumo como sinónimo de riqueza y bienestar. El mercado ha impulsado dinámicas de extracción de recursos irresponsables para crear bienes manufacturados. No extenderé las causas del consumismo ya que están bien delimitadas en la guía de comisión. 

\subsection{Consecuencias }

\noindent La extracción excesiva de recursos naturales ha causado diversos problemas ambientales, tanto del lado de la extracción de los materiales, como de los desechos generados cuando se termina la vida útil del producto. Todo esto hace parte del problema del cambio climático al que las naciones se están enfrentando —y que de  hecho se tratan en otros comités— tan desesperadamente.

Otra consecuencia que poco se discute pero que es igualmente relevante, son los desastres sociales de multinacionales mineras en países en vías de desarrollo. Se crean minas ilegales con condiciones laborales espantosas para los trabajadores de la comunidades aledañas.

\subsection{Políticas de la delegación frente al tema}

\section{Antecedentes }

\noindent Amazon es una empresa que conecta a vendedores de todo tipo de productos con los consumidores. Somos conscientes de que gran parte de los productos que están en nuestra plataforma fomentan el consumismo. Sin embargo, no nos hacemos responsables por los daños ambientales de los bienes que se venden en nuestra plataforma. Solo queremos hacer felices a nuestros clientes con el servicio, pero de ellos depende hacer una compra consciente. 

Ahora bien, nuestra empresa ha ido innovando los procesos de \textit{delivery} para reducir nuestra huella ambiental. En este momento usamos tanto automóbiles como aviones a base de gasolina para transportar la mayoría de productos; pero estamos invirtiendo grandes cantidades de dinero en tecnologías que no contaminen como robots eléctricos de entregas en las aceras y drones para entregar productos a larga distancia. Estamos a la espera del desarrollo de servicios de transporte que utilicen electricidad. Notoriamente, todavía no existen aviones eléctricos de carga. En cualquier caso, estamos dispuestos a colaborar con cualquier tecnología que beneficie al ambiente sin sacrificar la experiencia de nuestros clientes. 

\subsubsection{Posibles soluciones}

\noindent Estamos convencidos de que los cambios sustanciales en el ambiente dependen de los consumidores, y de su disposición para comprar productos que no contaminen, incluso aunque estos tengan un precio mayor. Por ende, impulsaremos iniciativas en el comité donde se conscientice a la ciudadanía la importancia de sus compras. 

Asimismo, apoyaremos propuestas que requieran el uso de tecnología o herramientas de análisis de datos. Tenemos la convicción de que la inteligencia artificial puede contribuir significantivamente a reducir los efectos del cambio climático.

\subsection{Preguntas Orientadoras }

\noindent Las preguntas orientadores se responden indirectamente en el resto del Folder.







































































\bibliography{Citas}




\end{document} %Acá termina
